\section{Seguridad de la política y regularización}

\subsection{Regularization para prevenir el colapso de la política}

Ya sea PPO o SAC, cuando la actualización de la política es demasiado agresiva, el desempeño puede colapsar. En la práctica, las técnicas de regularization más comunes incluyen:

\begin{itemize}
    \item \textbf{KL penalty}: se añade $\lambda D_{\mathrm{KL}}(\pi_{\theta'}\|\pi_\theta)$ al objective, penalizando directamente las desviaciones excesivas.
    \item \textbf{Entropy decay}: en SAC se ajusta $\alpha$ para que la política reduzca gradualmente la exploración hacia el final del entrenamiento y así converja.
    \item \textbf{Gradient clipping}: limita la norma del gradiente para evitar pasos demasiado grandes.
\end{itemize}

\subsection{Monitoreo del comportamiento de la política}

Métricas de monitoreo recomendadas:

\begin{itemize}
    \item \textbf{Average KL divergence} between successive policies
    \item \textbf{Episode variance} of returns—high variance signals instability
    \item \textbf{Q-value rollout test}—simulate fixed rollout to check escalation
\end{itemize}

\begin{knowledgebox}{Diseño de alertas automáticas}
Establecer umbrales (por ejemplo, generar un error cuando KL > 0.05, o detener el entrenamiento cuando la episode std dev aumenta 30\%) permite que el pipeline regrese automáticamente a un prior checkpoint. Este mecanismo se usa ampliamente en proyectos de Walls-Free RL (Sim2Real).
\end{knowledgebox}

\subsection{Resumen de la sección}

La programación segura requiere combinar regularization con métricas de monitoreo. KL penalty, entropy decay y gradient clipping son las técnicas comunes para evitar desviaciones excesivas, mientras que el monitoreo de KL, variance y Q-value rollout ofrece alertas de riesgo en tiempo real.

